\begin{textsample}{POS Dim 4 – 2023\_c5 – Score -1.00 – 2023\_c5/enade\_2023\_farm\_a515e48704.txt}  \label{ex:f4_pos_269}
CAPÍTULO 4 DISTRIBUIÇÃO DE CONCEITOS A finalidade deste capítulo é apresentar um panorama nacional sobre a distribuição do Conceito Enade 2023.
Para traçar esse panorama, recorre -se a tabelas e gráficos, que são acompanhados de comentários e análises que contribuem para a contextualização dos dados neles expostos.
Para visualizar o panorama do Enade 2023, em nível nacional, na primeira seção deste capítulo, os dados sobre o número de cursos participantes são organizados por grande região.
Na segunda, o mapeamento do conceito do exame é feito por categoria administrativa das IES e modalidade de oferta dos cursos, de maneira cruzada às grandes regiões.
Na última seção, a lógica de apresentação dos dados é feita por organização acadêmica e, também, por grande região de funcionamento dos cursos, no país como um todo. 4.1 PANORAMA NACIONAL DA DISTRIBUIÇÃO DOS CONCEITOS A lógica adotada na apresentação dos dados, nesta seção e nas duas seguintes, pressupõe o registro de duas informações de natureza reiterativa.
O primeiro diz respeito aos valores possíveis do Enade, que variam em níveis de 1 a 5, sendo que, à medida que esses valores aumentam, o curso é melhor avaliado, e ao fato de que aqueles cursos que tiveram somente um ou nenhum concluinte participante não obtêm o Conceito Enade, ficando sem conceito ( SC ), como ressaltado na apresentação deste relatório.
O segundo registro é que, caso haja diferença entre os cursos tabulados neste capítulo e no capítulo 1, ela corresponderá aos cursos que não foram avaliados, ou seja, àqueles, em princípio, sem estudantes concluintes que tivessem sido inscritos para a prova.
Na Tabela 4.1, são apresentadas a quantidade e a distribuição de cursos de Farmácia participantes do Enade 2023, por faixa de conceito e grande região.
A partir dos dados da Tabela 4.1, nota -se que, dos 590 cursos participantes, 225 ( 38,1 \% ) classificaram -se com conceito 2, o valor modal.
Esse foi também o conceito modal nas regiões : Centro-Oeste ( 40,0 \% ), Nordeste ( 52,1 \% ) e Norte ( 53,2 \% ).
O conceito 3 foi o segundo mais frequente em nível nacional ( 30,2 \%, correspondendo a 178 cursos ), tendo sido o segundo conceito mais frequente na maioria das regiões, exceto nas regiões Sudeste e Sul, em que foi o conceito modal.
O conceito 4 foi o terceiro ( 17,1 \%, correspondendo a 101 cursos ).
Houve, ainda, 38 ( 6,4 \% ) cursos que receberam o conceito 5 e 36 ( 6,1 \% ) cursos que receberam o conceito 1.
Dos 590 cursos de Farmácia, 12 ( 2,0 \% ) ficaram sem conceito ( SC ).
A região Norte participou com 47 cursos, o que corresponde a 8,0 \% do total nacional.
Desses, 25 cursos receberam o conceito 2, o que equivale a 53,2 \% do total regional, o conceito modal na região, conforme mencionado.
A oito cursos ( 17,0 \% ) atribuiu -se o conceito 3 ; a oito cursos ( 17,0 \% ), o conceito 1 ; e a cinco cursos ( 10,6 \% ) o conceito 4.
Ainda nessa região, nenhum curso ( 0,0 \% ) obteve conceito 5 e um curso ( 2,1 \% ) ficou sem conceito ( SC ).
A região Nordeste participou com 144 cursos, correspondentes a 24,4 \% do total nacional.
Nessa região, 75 cursos ( 52,1 \% em termos regionais ) obtiveram conceito 2, o conceito modal para a região, como já comentado.
Os conceitos 3 e 4 foram atribuídos, respectivamente, a 34 e a 16 cursos ( 23,6 \% e 11,1 \%, respectivamente ).
Ainda nessa região, oito cursos ( 5,6 \% ) foram avaliados com o conceito 5 e sete cursos ( 4,9 \% ) foram avaliados com o conceito 1.
Outros quatro cursos ( 2,8 \% ) ficaram sem conceito ( SC ).
Para os 230 cursos participantes da região Sudeste ( 39,0 \% do total nacional ), o conceito 3 foi o conceito modal, como já anteriormente assinalado, tendo sido obtido por 86 cursos ( 37,4 \% ).
O conceito 2 foi atribuído a 74 cursos ( 32,2 \% ) e o conceito 4, a 48 cursos ( 20,9 \% ).
Receberam os conceitos 5 e 1, respectivamente, 14 cursos ( 6,1 \% ) e cinco cursos ( 2,2 \% ).
Nessa região, três cursos ( 1,3 \% ) ficaram sem conceito ( SC ).
Os 99 cursos da região Sul corresponderam a 16,8 \% do total nacional.
Como já indicado, houve predominância do conceito 3 ( conceito modal ), atribuído a 32,3 \% dos cursos da região, ou seja, atribuído a 32 dos 99 cursos participantes na região Sul.
O conceito 4 foi atribuído a 23 cursos ( 23,2 \% ) ; o conceito 2, 23 cursos ( 23,2 \% ) ; e o conceito 5, a 13 cursos ( 13,1 \% ).
Na região Sul, sete cursos ( 7,1 \% ) obtiveram conceito 1 e um curso ( 1,0 \% ) ficou sem conceito ( SC ).
Já dos 70 cursos participantes na região Centro-Oeste ( 11,9 \% do total nacional ), 28 ( 40,0 \% em termos regionais ) receberam conceito 2, o conceito modal, como já destacado. 18 cursos ( 25,7 \% ) obtiveram o conceito 3, nove cursos ( 12,9 \% ) obtiveram o conceito 4 e nove cursos ( 12,9 \% ), o conceito 1.
Nessa região, três cursos ( 4,3 \% ) obtiveram o conceito 5 e três cursos ( 4,3 \% ) ficaram sem conceito ( SC ).
Em um gráfico de distribuição cumulativa, em cada valor de abscissa ( eixo x ), o valor apresentado no eixo das ordenadas ( eixo y ) é o valor acumulado da variável ( Conceito Enade ).
Em linhas gerais, em um gráfico desta mesma tipologia, poligonais mais à esquerda representam, no caso dos dados apresentados nesta seção, as grandes regiões com a pior distribuição, e poligonais mais à direita, as grandes regiões com os melhores conceitos.
A partir destas diretrizes, no Gráfico 4.1, é apresentada a distribuição cumulativa do Conceito Enade, segundo a grande região.
No Gráfico 4.1, observando a linha vertical pontilhada, que ajuda a visualizar o percentual de conceitos acumulados até o valor 2, ou seja, conceitos 1 e 2, pode -se notar, por exemplo, que a região Norte ( linha azul ) apresenta 71,7 \% dos cursos com conceito 2 ou abaixo deste valor.
A região Sul ( linha roxa ), por outro lado, apresenta 30,6 \% dos cursos na mesma situação : com conceito 2 ou abaixo.
Destaca -se que, a região Norte, o conceito 5 não foi alcançado por nenhum curso, na qual todos os cursos têm conceito menor ou igual a 4, e a respectiva poligonal cumulativa alcança 100 \% no conceito 4.
A região Norte apresenta os cursos com conceitos mais baixos, ao passo que a região Sul ( linha roxa ) apresenta a distribuição com valores maiores.
Os cursos em IES situadas na região Centro-Oeste ( linha laranja ) apresentam uma situação intermediária. 4.2 CONCEITOS POR CATEGORIA ADMINISTRATIVA, MODALIDADE DE OFERTA DOS CURSOS E GRANDE REGIÃO Na Tabela 4.2, é apresentada a distribuição dos cursos de Farmácia participantes do Enade 2023, por categoria administrativa e modalidade de oferta, de acordo com os conceitos por eles alcançados, segundo as grandes regiões brasileiras.
De acordo com as informações da Tabela 4.2, dos 590 cursos participantes, 71 ( 12,0 \% ) eram ministrados em instituições públicas e 519 ( 88,0 \% ), em instituições privadas.
Quanto à modalidade de oferta, 561 ( 95,1 \% ) cursos eram presenciais e 29 ( 4,9 \% ), a distância.
De acordo com as informações da mesma tabela, em termos nacionais, entre instituições públicas e privadas, observa -se que, dos 38 cursos avaliados com conceito 5, 23 foram oferecidos em IES públicas, e 15, em IES privadas.
Dos 519 cursos participantes de IES privadas, o conceito 2 foi o valor modal, atribuído a 222 ( 42,8 \% ) cursos.
Em relação aos demais cursos de instituições privadas, 172 cursos ( 33,1 \% ) foram avaliados com o conceito 3, 63 cursos ( 12,1 \% ) foram avaliados com o conceito 4, 35 cursos ( 6,7 \% ) foram avaliados com o conceito 1 e, 15 cursos ( 2,9 \% ), com o conceito 5.
Nessa categoria, 12 cursos ( 2,3 \% ) ficaram sem conceito ( SC ).
Na rede pública, o conceito modal foi o 4, atribuído a 38 cursos, correspondentes a 53,5 \% dos 71 cursos da categoria.
Dos demais cursos participantes, 23 cursos ( 32,4 \% ) obtiveram o conceito 5, seis cursos ( 8,5 \% ) receberam o conceito 3 e três cursos ( 4,2 \% ), o conceito 2.
O conceito 1 foi atribuído a um curso ( 1,4 \% ).
Nessa categoria administrativa, nenhum curso ( 0,0 \% ) ficou sem conceito ( SC ).
Dos 561 cursos participantes da modalidade presencial, o conceito 2 foi o valor modal, atribuído a 209 ( 37,3 \% ) cursos.
Dos demais cursos dessa modalidade de oferta, 175 cursos ( 31,2 \% ) receberam o conceito 3 ; 100 cursos ( 17,8 \% ) receberam o conceito 4 ; 38 cursos ( 6,8 \% ), o conceito 5 ; e 30 cursos ( 5,3 \% ), o conceito 1.
Nessa modalidade de oferta, nove cursos ( 1,6 \% ) ficaram sem conceito ( SC ).
Dos 29 cursos participantes da modalidade a distância, o conceito 2 foi o valor modal, atribuído a 16 ( 55,2 \% ) cursos.
Dos demais cursos ofertados nessa modalidade, seis cursos ( 20,7 \% ) receberam conceito 1, três cursos ( 10,3 \% ) recebeu o conceito 3 ; um curso ( 3,4 \% ) recebeu o conceito 4 ; e nenhum curso ( 0,0 \% ) recebeu conceito 5.
Nessa modalidade de oferta, três cursos ( 10,3 \% ) ficaram sem conceito ( SC ).
Na análise por região, observa -se, que, na região Norte, as instituições públicas participaram com seis cursos ( 12,8 \% total regional ), dos quais nenhum ficou sem conceito ( SC ) ou recebeu o conceito 5.
O conceito modal foi 4, atribuído a quatro cursos, correspondendo a 66,7 \% dessa categoria na região.
O conceito 3 teve uma concentração próxima da modal, com um curso ( 16,7 \% ), e o conceito 1 também teve um curso ( 16,7 \% ).
Nenhum curso recebeu os conceitos 5 e 2.
As instituições privadas participaram com 41 cursos na região Norte ( 87,2 \% do total regional ), e um curso ficou sem conceito ( SC ).
O conceito modal nas instituições privadas na região foi o 2, atribuído a 25 cursos, o que corresponde a 61,0 \% dessa categoria na região.
Os conceitos 3, 1 e 4 foram atribuídos a sete, sete e um curso, respectivamente.
O conceito 5 foi atribuído a nenhum curso ( 0,0 \% ).
Um curso ficou sem conceito ( SC ).
Dos cursos oferecidos na região Norte, 46 eram presenciais e um ficou sem conceito ( SC ).
Dos 46 cursos oferecidos, 25 obtiveram o conceito 2 ( modal ) ; oito, o conceito 3 ; sete, o conceito 1 ; cinco, o conceito 4 ; e nenhum, o conceito 5.
Um curso ficou sem conceito ( SC ).
Um curso era na modalidade a distância, tendo um recebido o conceito modal 1.
Na região Nordeste, a rede privada concentra 124 dos 144 cursos participantes, o equivalente a 86,1 \% do total da região.
Entre os cursos oferecidos por IES privadas no Nordeste, 73 cursos foram avaliados com conceito 2, o conceito modal.
Trinta e quatro cursos obtiveram o conceito 3 ; sete cursos, o conceito 1 ; cinco cursos, o conceito 4 ; e um curso obteve o conceito 5.
Nessa categoria administrativa e grande região, quatro cursos ficaram sem conceito ( SC ).
As instituições públicas da região Nordeste participaram com 20 cursos ( 13,9 \% do total da região ).
Desses, 11 ( 55,0 \% ) obtiveram conceito 4, o valor modal, seguidos por sete cursos, que receberam o conceito 5.
O conceito 2 foi atribuído a dois cursos ; os conceitos 1 e 3, a nenhum curso.
Nenhum curso ficou sem conceito ( SC ).
Dos cursos oferecidos na região Nordeste, 141 eram presenciais : três ficaram sem conceito ( SC ), sete obtiveram o conceito 1 ; 73, o conceito 2, o valor modal ; 34, o conceito 3 ; 16, o conceito 4 ; e oito, o conceito 5.
Três cursos eram na modalidade a distância, tendo dois cursos recebido o conceito modal 2 ; nenhum, o conceito 1 ; nenhum, o conceito 3 ; nenhum, o conceito 4 ; e nenhum, o conceito 5.
Um curso ficou sem conceito ( SC ).
Na região Sudeste, a proporção de cursos da rede privada, 90,4 \%, correspondeu a 208 dos 230 cursos participantes.
Nessa combinação de categoria administrativa e grande região, o conceito modal foi o 3, atribuído a 84 cursos ; e três cursos ficaram sem conceito ( SC ).
Os demais foram avaliados com os seguintes conceitos : 2 ( 74 cursos ), 4 ( 36 cursos ), 5 ( seis cursos ) e 1 ( cinco cursos ).
Dos 22 cursos oferecidos por instituições públicas na região Sudeste, oito obtiveram o conceito 5 ; 12, o conceito 4 ; dois, o conceito 3 e nenhum obteve o conceito 2 ou 1.
Nenhum curso ficou sem conceito ( SC ).
Ainda na região Sudeste, grande parcela dos cursos ( 94,8 \% ) foi oferecida na modalidade presencial : 218 dos 230 cursos.
Para esses, a moda foi o conceito 3, com 84 cursos.
Os demais cursos presenciais da região foram alocados nos conceitos 2 ( 66 cursos ), 4 ( 47 cursos ), 5 ( 14 cursos ) e 1 ( quatro cursos ).
Três cursos presenciais ficaram sem conceito ( SC ).
Na região Sudeste, 12 cursos foram oferecidos na modalidade a distância, tendo sido um com conceito 1 ; oito com conceito 2 ; dois com conceito 3 ; um com conceito 4 ; e nenhum com conceito 5.
Nenhum curso ficou sem conceito ( SC ).
As instituições privadas concentraram 85 dos 99 cursos participantes da região Sul, 85,9 \% do total regional.
Desses, 31 cursos obtiveram conceito 3, o conceito modal.
Nessa combinação de categoria administrativa e grande região, 22 receberam o conceito 2 ; 17 receberam o conceito 4 ; sete, o conceito 5 ; e sete obtiveram o conceito 1.
Um curso ficou sem conceito ( SC ).
As instituições públicas na região Sul participaram com 14 cursos ( 14,1 \% dos cursos da região ), aos quais foram atribuídos os conceitos 1 ( nenhum curso ), 2 ( um curso ), 3 ( um curso ), 4 ( seis cursos ) e 5 ( seis cursos ).
Nenhum curso da região Sul oferecido por IES públicas ficou sem conceito ( SC ).
Na região Sul, 87 dos cursos foram oferecidos na modalidade presencial.
Para esses, a moda foi o conceito 3, com 31 cursos.
Os demais cursos presenciais da região foram alocados nos conceitos 4 ( 23 cursos ), 2 ( 17 cursos ), 5 ( 13 cursos ) e 1 ( três cursos ).
Nenhum curso ficou sem conceito ( SC ).
Na modalidade a distância, foram oferecidos 12 cursos, sendo quatro cursos com conceito 1, seis cursos com conceito 2 ( modal ), um curso com conceito 3, nenhum curso com conceito 4 e nenhum curso com conceito 5.
Um curso ficou sem conceito ( SC ) nessa região.
Na região Centro-Oeste, 61 dos 70 cursos participantes eram de IES privadas ( 87,1 \% em termos regionais ).
Desses, 28 cursos receberam conceito 2, o valor modal.
Dos demais cursos, 16 receberam o conceito 3, nove receberam o conceito 1, quatro receberam o conceito 4 e um recebeu o conceito 5.
Três ficaram sem conceito ( SC ).
Quanto aos nove cursos oferecidos por instituições públicas na região Centro-Oeste ( 12,9 \% do total regional ), nenhum ficou com conceito 1 ; nenhum, com conceito 2 ; dois, com conceito 3 ; cinco, com conceito 4 ; e dois ficaram com conceito 5.
Nenhum ficou sem conceito ( SC ) nessa região.
Também na região Centro-Oeste, a maioria dos cursos foi ofertada na modalidade presencial : 69 ( 98,6 \% ).
Para esses, a moda foi o conceito 2, com 28 cursos ( 40,6 \% ).
Os demais cursos presenciais da região foram alocados nos conceitos 3 ( 18 cursos, 26,1 \% ), 4 ( nove cursos, 13,0 \% ), 1 ( nove cursos, 13,0 \% ) e 5 ( três cursos, 4,3 \% ).
Dois cursos presencial da região ficaram sem conceito ( SC ).
Na modalidade a distância, foi oferecido um curso.
O único curso a distância nessa região ficou sem conceito ( SC ).
A título de síntese, no Gráfico 4.2, apresenta -se a distribuição do Conceito Enade, segundo a categoria administrativa da IES.
Nesse gráfico, como dito anteriormente, poligonais mais à esquerda representam Conceito Enade mais baixos, e poligonais mais à direita, Conceito Enade mais altos.
É importante observar que, para cada poligonal presente nesse gráfico, a porcentagem constante no eixo y refere -se ao universo de cursos oferecidos por instituições públicas ( poligonal azul ) ou privadas ( poligonal verde ), respectivamente.
Os cursos em IES públicas ( poligonal azul ) apresentam uma poligonal mais à direita do que os em IES privadas ( poligonal verde ) e, consequentemente, uma distribuição de Conceito Enade com valores maiores. 4.3 CONCEITOS POR ORGANIZAÇÃO ACADÊMICA E GRANDE REGIÃO Na Tabela 4.3, encontra -se a distribuição dos conceitos atribuídos aos cursos participantes do Enade 2023, na área de Farmácia, por organização acadêmica, segundo as grandes regiões brasileiras.
Como mostram os dados da Tabela 4.3, dos 590 cursos de Farmácia participantes, 199 eram oferecidos em universidades, 215 em centros universitários e 174 em faculdades ; dois cursos foram oferecidos por CEFET/IF.
Essa distribuição corresponde a, respectivamente, 33,7 \%, 36,4 \%, 29,5 \% e 0,3 \% dos cursos.
De acordo com os dados apresentados na mesma tabela, de todos os 101 cursos avaliados com o conceito 4, 60 eram vinculados a universidades.
Os cursos vinculados a esse tipo de organização acadêmica tiveram o conceito 3 como conceito modal, atribuído a 63 cursos ( 31,7 \% ).
Os demais cursos ofertados por universidades receberam os conceitos 2 ( 39 cursos ), 5 ( 30 cursos ) e 1 ( seis cursos ).
Um curso ficou sem conceito ( SC ).
Dos cursos vinculados a centros universitários, o conceito modal foi o 2, atribuído a 92 cursos ( 42,8 \% ).
Seis cursos vinculados a esse tipo de organização acadêmica ficaram sem conceito ( SC ), e os outros cursos receberam os conceitos 3 ( 72 cursos ), 4 ( 25 cursos ), 1 ( 13 cursos ) e 5 ( sete cursos ).
Dos 174 cursos mantidos por faculdades, 93 ( 53,4 \% ) receberam conceito 2, o conceito modal.
O segundo conceito mais frequentemente atribuído aos cursos oferecidos por essa categoria de IES foi o conceito 3, recebido por 43 cursos.
Na sequência, 17 cursos obtiveram conceito 1 ; 15 cursos, o conceito 4 ; e um curso, o conceito 5.
Dos cursos oferecidos em faculdades, cinco cursos ficaram sem conceito ( SC ).
Dos cursos oferecidos por CEFET/IF, os conceitos modais foram o 2 e 4, atribuído a um curso ( 50,0 \% ) cada.
Nenhum curso vinculado a esse tipo de organização acadêmica ficou sem conceito ( SC ) ou recebeu os outros conceitos.
Considerando -se, separadamente, as regiões brasileiras, verifica -se que, na região Norte, as universidades participaram com dez cursos avaliados, aos quais foram atribuídos os seguintes conceitos : 1 a um curso ; 2 a quatro cursos ; 3 a um curso ; conceito 4, conceito modal, a quatro cursos ; e 5 a nenhum curso.
Nenhum curso ficou sem conceito ( SC ).
Os centros universitários da região Norte foram representados por 19 cursos, dos quais quatro cursos obteve conceito 1 ; oito cursos, o valor modal, conceito 2 ; cinco cursos, conceito 3 ; um curso, o conceito 4 ; e nenhum curso, o conceito 5.
Um curso ficou sem conceito ( SC ).
As faculdades participaram com 18 cursos na região Norte, os quais receberam os conceitos 1 ( três cursos ), 2 ( 13 cursos ), 3 ( dois cursos ), 4 ( nenhum curso ) e 5 ( nenhum curso ).
Nenhum curso ficou sem conceito ( SC ).
Os CEFET/IF não participaram com cursos.
Na região Nordeste, as universidades participaram com 27 dos 144 cursos da área de Farmácia oferecidos.
O conceito modal foi 4, atribuído a 12 cursos.
Outros sete cursos obtiveram o conceito 5 ; cinco cursos, o conceito 2 ; três cursos, o conceito 3 ; e nenhum curso, o conceito 1.
Nenhum curso ficou sem conceito ( SC ).
Os centros universitários contaram com 57 cursos participantes na região Nordeste, 37 dos quais receberam o conceito modal 2.
Para os demais cursos, foram atribuídos os conceitos 3 ( 16 cursos ), 1 ( dois cursos ), 5 ( um curso ) e 4 ( nenhum curso ).
Um curso oferecido por centros universitários no Nordeste ficou sem conceito ( SC ).
As faculdades foram representadas por 60 cursos na região Nordeste, os quais receberam os conceitos 1 ( cinco cursos ), 2 ( 33 cursos ), 3 ( 15 cursos ), 4 ( quatro cursos ) e 5 ( nenhum curso ).
Três cursos ficaram sem conceito ( SC ).
Na região Nordeste, os CEFET/IF não participaram com cursos.
Na região Sudeste, as universidades concentraram 92 dos 230 cursos de Farmácia da região.
Dos cursos oferecidos em IES com este tipo de organização acadêmica, nessa região, o conceito modal foi 3, atribuído a 36 cursos.
Os demais cursos receberam os conceitos 4 ( 23 cursos ), 2 ( 22 cursos ), 5 ( oito cursos ) e 1 ( dois cursos ).
Um curso ficou sem conceito ( SC ).
Os centros universitários participaram com 81 cursos na região Sudeste, dos quais 30 obtiveram o conceito modal 2.
O conceito 3 foi atribuído a 29 cursos ; o conceito 4 a 14 cursos ; o conceito 5 a cinco cursos ; e o conceito 1 a dois cursos.
Um curso ficou sem conceito ( SC ) nessa combinação de organização acadêmica e região.
Já as faculdades foram representadas por 56 cursos na região Sudeste, dos quais um ficou sem conceito ( SC ).
O conceito modal 2 foi atribuído a 22 cursos.
Aos demais cursos foram atribuídos os conceitos 3 ( 21 cursos ), 4 ( dez cursos ), 5 ( um curso ) e 1 ( um curso ).
Os CEFET/IF participaram com um curso, o qual recebeu conceito 4.
Dos 99 cursos da região Sul, 53 eram mantidos por universidades, para os quais o conceito modal foi 3, atribuído a 18 cursos.
Os demais cursos receberam os conceitos 4 ( 15 cursos ), 5 ( 12 cursos ), 2 ( seis cursos ) e 1 ( dois cursos ).
Em IES com esse tipo de organização acadêmica, nenhum curso da região Sul ficou sem conceito ( SC ).
Dos 33 cursos participantes de centros universitários da região Sul, um ficou sem conceito ( SC ), e 12 cursos receberam conceito 3, o valor modal.
Os conceitos restantes ficaram assim distribuídos : 2 ( nove cursos ), 4 ( sete cursos ), 1 ( três cursos ) e 5 ( um curso ).
Quanto aos 12 cursos vinculados a faculdades na região Sul, sete cursos receberam o conceito 2, valor modal.
Os demais conceitos ficaram assim distribuídos : 3 ( dois cursos ), 1 ( dois cursos ), 4 ( um curso ) e 5 ( nenhum curso ).
Dos 12 cursos, nenhum ficou sem conceito ( SC ).
Os CEFET/IF participaram com um curso na região Sul, ao qual foi atribuído o conceito 2.
Na região Centro-Oeste, 17 dos 70 cursos oferecidos eram mantidos por universidades.
Para esse tipo de organização acadêmica, o conceito modal foi 4, atribuído a seis cursos.
Em relação às universidades, nenhum curso da região ficou sem conceito ( SC ).
Os demais cursos receberam os conceitos 3 ( cinco cursos ), 5 ( três cursos ), 2 ( dois cursos ) e 1 ( um curso ).
Os centros universitários da região Centro-Oeste contaram com 25 cursos : dez cursos foram avaliados com o conceito 3 ( modal ) ; oito cursos, com o conceito 2 ; três cursos, com o conceito 4 ; dois cursos, com o conceito 1 ; e nenhum curso obteve o conceito 5.
Dois cursos ficaram sem conceito ( SC ).
Dos 28 cursos oferecidos em faculdades na região Centro-Oeste, 18 cursos receberam o conceito 2 ( modal ).
Os conceitos restantes ficaram assim distribuídos : 1 ( seis cursos ), 3 ( três cursos ), 5 ( nenhum curso ) e 4 ( nenhum curso ).
Nessa região, um curso oferecido em faculdades ficou sem conceito ( SC ).
No que se refere aos CEFET/IF da região Centro-Oeste, nenhum curso foi avaliado.
De maneira semelhante à forma como se procedeu em relação à seção anterior, esta seção encerra -se com o Gráfico 4.3, que apresenta a distribuição do Conceito Enade dos cursos de Farmácia segundo a organização acadêmica da IES.
No Gráfico 4.3, apresenta -se a distribuição do Conceito Enade segundo a organização acadêmica da IES.
A leitura desse último gráfico é semelhante à dos anteriores, seguindo -se a lógica de que poligonais deslocadas para a esquerda correspondem aos conceitos mais baixos ( 1, 2 e 3 ) e poligonais deslocadas para a direita correspondem a Conceitos Enade mais altos.
Nas universidades ( linha azul ), todos os cursos obtiveram conceito 5 ou abaixo disso, já que a poligonal cumulativa alcança 100 \% para o conceito 5.
Além disso, os cursos de universidades apresentam uma poligonal mais à direita do que as demais, denotando uma melhor distribuição de Conceitos Enade.
No outro extremo, encontram -se os cursos mantidos por faculdades ( linha vermelha ), apresentando as respectivas poligonais mais à esquerda e, consequentemente, uma distribuição de Conceitos Enade com os menores valores.
Os cursos mantidos por centros universitários ( linha verde ) apresentam poligonais mais centrais e, consequentemente, uma distribuição de Conceitos Enade com os valores intermediários.

% matched lemmas: 
\end{textsample}
